\section{Discussion and Limitations}
\label{sec:discuss}

OrbitLLM's main point is architectural: LEO/NTN mobile-edge schedulers need a runtime interface that exposes the quantities binding large-model inference. Bits and CPU cycles are insufficient when memory fit depends on KV cache length and energy depends on the prefill/decode mix. Making these quantities explicit also makes the framework reproducible: changing hardware, quantization, or model family only changes the profile CSV, not the optimizer or simulator.

The most important limitation is that the reported numbers are RTX 4090 availability-anchor measurements, not spacecraft accelerator measurements. The hardware-scaling experiment is intended to test sensitivity to this choice, not to validate a particular on-orbit device. Radiation effects, spacecraft integration, accuracy changes under quantization, and detailed power replenishment are outside the scope of this workshop version.

PRE is also an abstraction rather than a single deployed vision module. We model it as lightweight semantic compression with retained-bit and utility-retention parameters, then sweep those parameters to expose when the action helps. A deployment would calibrate this curve with a concrete ROI selector, detector, or small vision-language encoder. Finally, our 24-hour traces are generated by a discrete-event simulator; future work should add real contact traces, multi-satellite interactions, and measured profiles from flight-relevant edge accelerators.

\section{Conclusion}
\label{sec:conclusion}
We presented OrbitLLM, a profile-anchored scheduling framework for large-model inference in LEO/NTN mobile edge systems. By turning LLM-specific inference behavior into a replaceable cost profile, the framework avoids the generic bits/cycles abstraction used by conventional offloading models. A simple value-density heuristic, calibrated against a MILP upper bound, outperforms All-Downlink, All-On-Board, and fixed-threshold baselines while respecting energy and memory constraints. Network, PRE, and hardware-scaling sweeps show how the ON/PRE/DOWN action mix adapts as mobile-edge resources change.

\section*{Reproducibility Statement}
The artifact contains the benchmark CLI, simulator, MILP bound, plotting scripts, configuration file, generated CSV results, and LaTeX table macros. The code and results are available at \url{https://github.com/1740919441a/OrbitLLM}. Running \texttt{python -m orbitllm.cli simulate} followed by \texttt{python -m orbitllm.cli plot} regenerates the reported tables and figures.
